\documentclass[11pt, letterpaper]{article}

% ── Packages ──────────────────────────────────────────────────────────────────
\usepackage[margin=1in]{geometry}
\usepackage{amsmath, amssymb}
\usepackage{graphicx}
\usepackage{booktabs}
\usepackage{hyperref}
\usepackage{enumitem}
\usepackage{titlesec}
\usepackage{parskip}
\usepackage{xcolor}
\usepackage{fancyhdr}
\usepackage{microtype}
\usepackage{cite}
\usepackage{array}
\usepackage{tabularx}

% ── Colours ───────────────────────────────────────────────────────────────────
\definecolor{accent}{RGB}{0, 70, 127}   % deep navy

% ── Hyperref ──────────────────────────────────────────────────────────────────
\hypersetup{
  colorlinks = true,
  linkcolor  = accent,
  citecolor  = accent,
  urlcolor   = accent,
}

% ── Section formatting ────────────────────────────────────────────────────────
\titleformat{\section}
  {\large\bfseries\color{accent}}{\thesection.}{0.5em}{}[\titlerule]
\titleformat{\subsection}
  {\normalsize\bfseries\color{accent!80!black}}{\thesubsection.}{0.5em}{}

% ── Header / Footer ───────────────────────────────────────────────────────────
\pagestyle{fancy}
\fancyhf{}
\fancyhead[L]{\small\textcolor{accent}{\textbf{HFT Order Intelligence}}}
\fancyhead[R]{\small\textcolor{gray}{Project Proposal}}
\fancyfoot[C]{\small\thepage}
\renewcommand{\headrulewidth}{0.4pt}

% ─────────────────────────────────────────────────────────────────────────────
\begin{document}

% ── Title block ───────────────────────────────────────────────────────────────
\begin{center}
  {\LARGE\bfseries\color{accent}
    Reverse-Engineering HFT Strategy from Reported Orders:\\[4pt]
    PnL Attribution, News-Event Reaction, and Feature Inference
  }\\[10pt]
  {\large Research Proposal}\\[4pt]
  {\normalsize \today}
\end{center}

\vspace{6pt}
\hrule height 1.2pt
\vspace{14pt}

% ── Abstract ──────────────────────────────────────────────────────────────────
\begin{center}
\begin{minipage}{0.92\textwidth}
\small\itshape
\textbf{Abstract.}
High-frequency trading (HFT) firms are required to report a subset of their
orders through regulatory and exchange mechanisms (FINRA OATS/CAT, 13-F filings,
and similar sources), yet these filings represent a sparse, time-delayed, and
often aggregated window into activity that occurs at microsecond resolution.
This project uses that sparse signal to reconstruct, as completely as possible,
the PnL profile of HFT-attributable orders, measure their latency response to
macroeconomic and firm-specific news events, and infer the likely feature set
and strategy class underpinning their trading behavior. The ultimate goal is to
build an open analytical framework that quantifies how much information is
recoverable about a market participant from their minimal regulatory footprint.
\end{minipage}
\end{center}

\vspace{10pt}

% ─────────────────────────────────────────────────────────────────────────────
\section{Motivation and Research Questions}
% ─────────────────────────────────────────────────────────────────────────────

HFT firms collectively account for a substantial fraction of U.S. equity volume
yet disclose only the minimum required by regulation. Each reported order is
therefore a data point of outsized informational value: it reveals a firm's
choice of instrument, direction, size, and timing under the specific market
conditions that prevailed at execution. Taken together, even a handful of
reported orders can constrain the strategy space considerably.

This project is organized around three core research questions:

\begin{enumerate}[label=\textbf{RQ\arabic*.}, leftmargin=2.2cm]
  \item \textbf{PnL Attribution.}
    Given a set of reported HFT orders and the corresponding limit-order-book
    (LOB) and trade-tape data, what is the realized PnL of those orders at
    various horizons (intraday mark-to-market, close-of-day, and multi-day
    holding if applicable)?

  \item \textbf{News-Event Reaction Timing.}
    How quickly do reported HFT orders respond to identifiable news events
    (FOMC announcements, earnings releases, macro data prints, and
    Twitter/news-sentiment spikes)? Does the reaction-time distribution imply
    co-location, direct feeds, or alternative data usage?

  \item \textbf{Strategy and Feature Inference.}
    Using the temporal, directional, and size patterns of reported orders, can
    we classify the likely strategy archetype (market-making, statistical
    arbitrage, momentum, latency arbitrage) and infer which features—spread,
    order imbalance, mid-price momentum, volatility regime—best explain the
    observed order placement behavior?
\end{enumerate}

% ─────────────────────────────────────────────────────────────────────────────
\section{Data Sources}
% ─────────────────────────────────────────────────────────────────────────────

\subsection{Reported Order Data}

\begin{itemize}[noitemsep]
  \item \textbf{CAT / OATS filings}: Consolidated Audit Trail reports providing
    NMS order-level data attributed to member firms, accessible via academic
    data agreements or FINRA research programs.
  \item \textbf{13-F quarterly holdings}: For HFT-affiliated entities that cross
    AUM thresholds, institutional holdings filings offer a coarse long-side
    snapshot.
  \item \textbf{Exchange proprietary feeds} (TAQ, ITCH): Used to reconstruct
    surrounding market context and to match reported orders against the public
    record.
\end{itemize}

\subsection{Market Microstructure Data}

\begin{itemize}[noitemsep]
  \item \textbf{LOBSTER / NASDAQ ITCH}: Full limit-order-book reconstructions
    at message-level granularity for selected symbols and dates.
  \item \textbf{NYSE TAQ}: Millisecond-stamped trade and quote data for broad
    cross-sectional coverage.
  \item \textbf{CME Globex MDP 3.0}: For futures instruments where HFT activity
    is pronounced (ES, NQ, ZN).
\end{itemize}

\subsection{News and Macro Event Data}

\begin{itemize}[noitemsep]
  \item \textbf{Refinitiv/LSEG NewsScope}: Timestamped news events with
    entity-level tagging.
  \item \textbf{FRED / BLS release calendar}: Macro data print timestamps
    (CPI, NFP, FOMC minutes) with pre-announced release times.
  \item \textbf{Ravenpack / Benzinga}: Alternative sentiment scores for
    real-time news impact estimation.
\end{itemize}

% ─────────────────────────────────────────────────────────────────────────────
\section{Methodology}
% ─────────────────────────────────────────────────────────────────────────────

\subsection{Step 1 — Order Matching and Context Reconstruction}

Each reported order $o_i = (t_i, s_i, d_i, q_i, p_i)$ (timestamp, symbol,
direction, quantity, price) is matched against the contemporaneous LOB snapshot
to recover:
\begin{itemize}[noitemsep]
  \item Best bid/ask and spread at time $t_i$
  \item Queue depth at the reported price level
  \item Estimated fill probability and expected slippage
\end{itemize}

For orders with no explicit fill confirmation in the public record, we use a
probabilistic fill model calibrated on execution-quality statistics from MIDAS
data.

\subsection{Step 2 — PnL Attribution}

Realized PnL for each filled order is computed at three horizons:

\begin{align}
  \text{PnL}^{(T)}_i &= d_i \cdot \bigl(m_{t_i + T} - p_i\bigr) \cdot q_i
\end{align}

where $m_{t_i + T}$ is the mid-price at horizon $T \in \{0.01\text{s},\;
0.5\text{s},\; 1\text{s},\; \text{EOD}\}$ and $d_i \in \{+1, -1\}$ encodes
buy/sell direction.

Aggregate PnL curves are constructed and bootstrapped confidence bands are
applied to assess statistical significance given the small reported-order sample.

\subsection{Step 3 — News-Event Reaction Analysis}

For each identified macro/news event $e_j$ with timestamp $\tau_j$:

\begin{enumerate}[noitemsep]
  \item Compute the set of HFT orders placed within window
    $[\tau_j - 5\text{s},\; \tau_j + 30\text{s}]$.
  \item Measure \textbf{reaction latency} $\Delta_j = t^*_j - \tau_j$, where
    $t^*_j$ is the first order placed after the event.
  \item Compare $\Delta_j$ against the distribution of all-participant first
    responses to the same event (from TAQ).
  \item Compute \textbf{directional accuracy}: fraction of post-event orders
    that correctly anticipate the subsequent price move.
\end{enumerate}

We additionally examine the \textbf{pre-event positioning}: order flow imbalance
in $[\tau_j - 60\text{s},\; \tau_j)$ attributed to the HFT entity, which may
reveal whether the firm trades on advance signal or purely on speed.

\subsection{Step 4 — Feature and Strategy Inference}

We treat the ordered sequence of $N$ reported orders as observations from a
latent decision process and attempt to infer the most likely feature set via
three complementary approaches:

\paragraph{4a. Descriptive signature analysis.}
Compute summary statistics (order-to-trade ratio, average holding time, size
clustering, time-of-day seasonality, spread-crossing frequency) and compare
against known benchmarks for each strategy archetype.

\paragraph{4b. Regression-based feature attribution.}
Regress order direction $d_i$ on contemporaneous LOB features:
\begin{align}
  \Pr(d_i = +1) &= \sigma\!\Bigl(\beta_0 + \beta_1 \cdot \text{OIB}_{t_i}
    + \beta_2 \cdot \text{Mom}_{t_i}^{(k)} + \beta_3 \cdot \hat{\sigma}_{t_i}
    + \beta_4 \cdot \text{Spread}_{t_i} + \varepsilon_i\Bigr)
\end{align}
where OIB is the LOB order-flow imbalance, Mom$^{(k)}$ is $k$-second
mid-price momentum, $\hat{\sigma}$ is realized volatility, and $\sigma(\cdot)$
is the logistic function. Significant $\beta$ coefficients identify which
microstructure features explain order direction.

\paragraph{4c. Classification against strategy templates.}
Build a labeled dataset of synthetic HFT order sequences generated from known
strategy archetypes (market-making, momentum, stat-arb, latency arb) using
simulation on real LOB data. Train a classifier (XGBoost or random forest)
and apply it to the reported orders.

% ─────────────────────────────────────────────────────────────────────────────
\section{Expected Outputs and Deliverables}
% ─────────────────────────────────────────────────────────────────────────────

\begin{tabularx}{\textwidth}{@{} l X @{}}
  \toprule
  \textbf{Deliverable} & \textbf{Description} \\
  \midrule
  PnL dashboard & Interactive visualization of per-order and aggregate PnL
    curves across horizons, with breakdown by symbol, date, and event type. \\[4pt]
  Reaction-time report & Ranked latency table of the HFT entity vs.\ market
    median across event categories; directional accuracy scorecard. \\[4pt]
  Feature attribution table & Regression output with coefficient estimates and
    significance levels for each candidate microstructure feature. \\[4pt]
  Strategy classification memo & Model output, confidence scores, and
    qualitative interpretation of the inferred strategy archetype. \\[4pt]
  Open-source codebase & Python package (pandas / polars / scikit-learn)
    covering data ingestion, LOB reconstruction, PnL engine, and
    classification pipeline; reproducible via GitHub. \\
  \bottomrule
\end{tabularx}


% ─────────────────────────────────────────────────────────────────────────────
\section{Related Work}
% ─────────────────────────────────────────────────────────────────────────────

This project builds on a body of literature in market microstructure and HFT
analysis:

\begin{itemize}[noitemsep]
  \item Brogaard, Hendershott, \& Riordan (2014) identify HFT order flow in
    NASDAQ data and document their superior ability to predict short-term price
    movements.
  \item Menkveld (2013) characterizes modern market-maker PnL decomposition
    into spread revenue and adverse-selection loss.
  \item Cont, Kukanov, \& Stoikov (2014) introduce the order-flow imbalance
    measure as a leading indicator of short-horizon price moves—directly
    applicable to our feature-attribution step.
  \item Baron, Brogaard, Hagströmer, \& Kirilenko (2019) document the
    heterogeneity of HFT strategies and their differential contribution to
    price discovery vs.\ intermediation.
\end{itemize}

% ─────────────────────────────────────────────────────────────────────────────
\section{Conclusion}
% ─────────────────────────────────────────────────────────────────────────────

Reported HFT orders are a rare and underexploited data source. Even their
sparse signal—timing, direction, size, and instrument—contains enough
information to reconstruct meaningful PnL estimates, characterize a firm's
responsiveness to information events, and constrain the space of plausible
trading strategies. This project provides a rigorous, reproducible framework
for that analysis and contributes an open-source toolkit to the market
microstructure research community. The methods developed here are also directly
applicable to regulatory transparency efforts and to academic work on
information asymmetry in modern equity markets.

\end{document}
